\input{preamble}

\begin{document}

\title{Exact Marginal--Regret Frontiers\\for Signed-Prefix Compression}

\author{\IEEEauthorblockN{Vu Phan An Nguyen}
\IEEEauthorblockA{\textit{Rider University} \\
Lawrenceville, NJ, USA}}

\maketitle

\begin{abstract}
A compression is often served as a removal order: a consumer who must shed $k$ items drops
the first $k$ and keeps the rest, and $k$ is not known in advance. Suppose, in addition, that
every item must be dropped about equally often at every budget. What does that cost, and how
does the answer depend on when the budget is chosen and on how exactly the weights are known?
We answer all three questions exactly for the signed-prefix family. There, dropping a set $S$
costs $|\sum_{i\in S}g_{i}|$ above the best achievable at that budget.

\emph{The frontier.} Widen the marginal band to $\eta/m$. The optimum $V_{m}(\eta)$ is then
known in closed form at every size and every $\eta\ge0$. For even $m$ it is
$\max\{\frac{2}{(m+4)(m-1)},\frac{4(1-\eta)}{m(m+4)}\}$. For odd $m$ it is the maximum of
two lines that cross at $\eta_{t}=1-\frac{2}{m(m-1)}$. It starts at $\frac{m+3}{m(m+4)}$ and
reaches the free value $\frac{2}{(m+4)(m-1)}$ exactly at $\eta=1$. So exact uniformity costs
a factor below $2$ in even size and a factor linear in $m$ in odd size.

\emph{A polynomial program.} The game over all $m!$ orders has exactly the value of a flow
over $O(m^{2})$ prefix compositions, with every per-item band carried.

\emph{Robustness.} For arbitrary weights the optimum is Lipschitz in $\|g-g'\|_{1}$. Perturb
by $\ell_{1}$ size $\varepsilon$, in odd size. While $\varepsilon=O(1/m)$ the price ratio is
linear in $m$, and beyond that it decays as $\Theta(1/\varepsilon)$. It is at least
$\frac{m-1}{4}$ below $\frac{1}{2m}$ and at least $\frac{1}{16\varepsilon}$ on
$[\frac1m,\frac16]$. Moving one light item by $\varepsilon$ gives at most
$\frac{10}{\varepsilon}$.

\emph{Timing.} Suppose the budget is chosen after the order is drawn. The free value is then
$\frac{2}{m+4}$, a factor $m-1$ larger. Under the band it is $\frac{2}{m+4}$ for even $m$
and $\frac{2+2(1-\eta)_{+}/m}{m+4}$ for odd $m$. So ex post the band costs at most a factor
$1+\frac1m$, and the two linear factors do not multiply.

Every theorem is machine-checked in Lean~4 at the paper's own objects.
\end{abstract}

\begin{IEEEkeywords}
ranking under constraints, prefix regret, linear programming duality, formal verification
\end{IEEEkeywords}

\section{Introduction}
\label{sec:intro}

\subsection{The question}

A served compression must often work at a size that is not known in advance. One way to
serve it is as a \emph{removal order}. The system returns an order of the items, and a
consumer who must shed $k$ of them drops the first $k$ and keeps the other $m-k$. Fairness or
exposure requirements then constrain the distribution of the order: each item should be
dropped about equally often at every budget. We ask what that constraint costs when the order
must be good at every budget simultaneously.

Fix weights $g=(g_{1},\dots,g_{m})$ and a normaliser $D>0$. For an order $\pi$ and a budget
$k$, the best achievable loss and the regret are
\begin{equation}
b_{k}(g)=\min_{|S|=k}\Bigl|\sum_{i\in S}g_{i}\Bigr|,\qquad
r_{k}(\pi;g)=\frac{\bigl|\sum_{\ell\le k}g_{\pi_{\ell}}\bigr|-b_{k}(g)}{D}.
\label{eq:regret}
\end{equation}
Here $\{\pi_{1},\dots,\pi_{k}\}$ is the \emph{dropped} set. In the physical model of the
compression, keeping the complement of a set $S$ costs $|\sum_{i\in S}g_{i}|/2$. So with
$D=\|g\|_{1}$, the regret~\eqref{eq:regret} is the loss of the compression that keeps the
last $m-k$ items, above the best compression of that size, normalised by $\|g\|_{1}/2$.

\leanbox{\lean{physical\_radius}, \lean{physical\_minimum}, \lean{physical\_normalized\_regret}
in \lean{LeanPaperCStabilityEndgameV1}}

For a distribution $P$ over orders, the constrained optimum is
\begin{equation}
V_{g}(\eta)=\inf_{P}\ \max_{0<k<m}\ \mathbb{E}_{\pi\sim P}\,r_{k}(\pi;g)
\label{eq:value}
\end{equation}
subject to the marginal band
\begin{equation}
\Bigl|\Pr\nolimits_{P}\bigl(i\in\{\pi_{1},\dots,\pi_{k}\}\bigr)-\frac{k}{m}\Bigr|\le\frac{\eta}{m}
\quad(\text{all }i,\ k\le m).
\label{eq:band}
\end{equation}
The band reads the same whether it is stated for dropping or for keeping. An item is kept at
size $m-k$ with probability within $\eta/m$ of $\frac{m-k}{m}$ exactly
when~\eqref{eq:band} holds. The budget in~\eqref{eq:value} is fixed before the draw, so the
objective is $\max_{k}\mathbb{E}\,r_{k}$. Section~\ref{sec:expost} treats the objective
$\mathbb{E}\max_{k}r_{k}$, where the budget is chosen after the draw. The two objectives are
never interchanged.

\subsection{Results}

We work on the family $g_{m}$ of~\eqref{eq:family} below. We write
$V_{m}(\eta)=V_{g_{m}}(\eta)$ and $V^{\mathrm{free}}_{m}=\frac{2}{(m+4)(m-1)}$. The latter is
the unconstrained optimum, because the band is vacuous for $\eta\ge m$.

\begin{enumerate}
\item[1)] \textbf{The whole frontier} (Theorems~\ref{thm:zero}, \ref{thm:even},
\ref{thm:odd}). $V_{m}(\eta)$ in closed form for every $m\ge3$ and every $\eta\ge0$. At
$\eta=0$ the price $V_{m}(0)/V^{\mathrm{free}}_{m}$ is $\frac{2(m-1)}{m}<2$ in even size and
$\frac{(m+3)(m-1)}{2m}$ in odd size. It vanishes exactly at $\eta_{*}=\frac{m-2}{2(m-1)}$
(even) and at $\eta_{*}=1$ (odd).
\item[2)] \textbf{A polynomial-size program with the same value} (Theorem~\ref{thm:flow}).
\item[3)] \textbf{Robustness, and where it ends} (Theorems~\ref{thm:robust},
\ref{thm:collapse}). A Lipschitz envelope for arbitrary weights, and the collapse scale of
the price. The price is linear exactly while the perturbation is $O(1/m)$, and it decays as
$\Theta(1/\varepsilon)$ beyond.
\item[4)] \textbf{Timing} (Proposition~\ref{prop:expost}, Theorem~\ref{thm:expostband}).
Choosing the budget after the draw multiplies the free value by $m-1$. Under the band the
ex-post value is known exactly, and there the band costs at most a factor $1+\frac1m$.
\end{enumerate}

\subsection{Prior work, and what does not follow from it}

Randomised rankings with constrained position or exposure marginals are a studied object.
Singh and Joachims~\cite{singh2018} introduce exposure constraints and optimise a linear
utility over the Birkhoff polytope. Celis, Straszak and Vishnoi~\cite{celis2018} bound group
representation in every prefix. Singh, Kempe and Joachims~\cite{singh2021} contrast ex-ante
and ex-post guarantees, as Freeman, Shah and Vaish~\cite{freeman2020} do in allocation. We
take that framing as given. We do not claim the flow technique either: reducing a permutation
polytope to a compositions DAG is standard~\cite{schrijver2003}.

What is new is the law that the reduction lets us prove, for an objective these works do not
treat. The signed-prefix regret is an absolute sum net of the per-budget optimum, taken at
every budget. It is not linear in the order, and the parity of $k$ enters through $b_{k}$.

\begin{table}[h]
\caption{Nearest sources, and a theorem of this paper that does not follow from each.}
\label{tab:prior}
\centering
\footnotesize
\begin{tabular}{@{}p{0.30\columnwidth}p{0.62\columnwidth}@{}}
\toprule
Nearest source & Theorem here that does not follow \\
\midrule
Exposure-constrained rankings~\cite{singh2018} & The exact signed-prefix frontier at every
size and band, with a $\Theta(m)$ parity split (Thms.~\ref{thm:zero}--\ref{thm:odd}) \\
Prefix-fair ranking~\cite{celis2018} & The exact breakpoints and the crossing point
$\eta_{t}$ of the odd frontier (Thm.~\ref{thm:odd}) \\
Ex-ante vs.\ ex-post~\cite{singh2021,freeman2020} & The exact banded ex-post value, and the
fact that the band and timing prices do not multiply (Thm.~\ref{thm:expostband}) \\
Flow reductions~\cite{schrijver2003} & Value equality with every per-item band carried
through relabelling (Thm.~\ref{thm:flow}) \\
\bottomrule
\end{tabular}
\end{table}

\subsection{What is proved, and how}

Three evidence classes appear, and they are never interchanged.
\begin{itemize}
\item \PROVED: a theorem stated at the paper's objects and checked in
Lean~4~\cite{demoura2021} against Mathlib~\cite{mathlib2020}.
\item \CERT: an exact rational certificate verified by a script.
\item \NUM: a floating-point solve or search.
\end{itemize}
Each theorem carries one line naming its Lean declaration. The proof architecture and every
supporting declaration are in the supplement. All modules compile with no \texttt{sorry} and
no \texttt{native\_decide}, and their axioms lie within \{\texttt{propext},
\texttt{Classical.choice}, \texttt{Quot.sound}\}. A compiled theorem certifies its statement
under its hypotheses. Every hypothesis of a headline is either a model assumption displayed
in it or discharged.

\section{Setting}
\label{sec:setting}

The family studied throughout is
\begin{equation}
g_{m}=\bigl(-3,\underbrace{-1,\dots,-1}_{p},\underbrace{1,\dots,1}_{q},3\bigr),
\qquad D=m+4,
\label{eq:family}
\end{equation}
with $p=\lceil\frac{m-2}{2}\rceil$ and $q=\lfloor\frac{m-2}{2}\rfloor$. The total is $0$ for
even $m$ and $-1$ for odd $m$, and $\|g_{m}\|_{1}=m+4=D$. Every weight is an odd integer, so
a sum of $k$ weights has the parity of $k$ and has absolute value at least $k\bmod2$. For
$g_{m}$ the per-budget optimum is exactly that,
\begin{equation}
b_{k}(g_{m})=k\bmod 2\qquad(0<k<m),
\label{eq:baseline}
\end{equation}
attained by an explicit subset at each budget. So on $g_{m}$ the physical regret, with the
baseline computed from its definition, is the normalised parity excess.

\leanbox{\lean{familyBaselineIdentification\_holds} in \lean{LeanPaperCBaselineAttainV1};
\lean{trueValue\_family} in \lean{LeanPaperCStabilityFamilyV1}}

The parity split lives on this structure. There are two heavy items of weight $\pm3$, and
light items of weight $\pm1$ whose counts differ by one in odd size. A $+3$ step between two
budgets of zero excess is impossible, because a zero-excess prefix sum has absolute value at
most $1$. That observation drives every odd lower bound below.

\section{The frontier}
\label{sec:frontier}

\begin{theorem}[Zero band]
\label{thm:zero}
For every even $m\ge4$ and every odd $m\ge3$,
\[
V_{m}(0)=
\begin{cases}
\dfrac{4}{m(m+4)}, & m\text{ even},\\[8pt]
\dfrac{m+3}{m(m+4)}, & m\text{ odd},
\end{cases}
\]
attained by an explicit law on all $m!$ orders of $g_{m}$. Hence
$V_{m}(0)/V^{\mathrm{free}}_{m}$ is $\frac{2(m-1)}{m}$ for even $m$ and
$\frac{(m+3)(m-1)}{2m}$ for odd $m$.
\end{theorem}

\leanbox{\lean{even\_value\_zero}, \lean{odd\_value\_zero}, \lean{odd\_uniform\_price\_ratio}
in \lean{LeanPaperCEndgameV1}}

\begin{theorem}[The even frontier]
\label{thm:even}
For every even $m\ge4$ and every $\eta\ge0$,
\[
V_{m}(\eta)=\max\Bigl\{\frac{2}{(m+4)(m-1)},\ \frac{4(1-\eta)}{m(m+4)}\Bigr\},
\]
so the price vanishes exactly at $\eta_{*}=\frac{m-2}{2(m-1)}$.
\end{theorem}

\leanbox{\lean{even\_frontier\_all} in \lean{LeanPaperCEndpointLawV1}}

\begin{theorem}[The odd frontier]
\label{thm:odd}
For every odd $m\ge3$ and every $\eta\ge0$,
\begin{multline*}
V_{m}(\eta)=\max\Bigl\{\frac{m+3-(m+1)\eta}{m(m+4)},\\
\frac{2}{(m+4)(m-1)}+\frac{(1-\eta)_{+}}{m(m+4)}\Bigr\}.
\end{multline*}
The two lines cross at $\eta_{t}=1-\frac{2}{m(m-1)}$, where
$V_{m}(\eta_{t})=\frac{2(m^{2}+1)}{m^{2}(m-1)(m+4)}$. In particular
$V_{m}(\eta)=V^{\mathrm{free}}_{m}$ exactly when $\eta\ge1$, and for any fixed $\eta<1$ the
odd price is $\Omega(m)$.
\end{theorem}

\leanbox{\lean{odd\_frontier}, \lean{odd\_kink\_value} in \lean{LeanPaperCOddFrontierV1};
\lean{odd\_breakpoint} in \lean{LeanPaperCEndpointLawV1}}

\emph{Why the frontier has this shape.} Every lower side comes from one budget or from an
average of budgets. Every upper side is one explicit law.

\emph{Zero band.} For even $m$ the weights sum to zero, so the prefix of length $m-1$ is minus
the last item. Then $r_{1}+r_{m-1}$ equals $\frac{2}{D}$ times the number of heavy items at
the two ends, and each end is heavy with probability $\frac{2}{m}$. For odd $m$ the prefix
of length $m-1$ is $-1$ minus the last item, and $m-1$ is even, so an item of weight $w$ in
last position costs $\frac{|1+w|}{D}$. Averaging over uniform last positions gives
$\frac{m+3}{mD}$. The upper laws are the following, each relabelled uniformly within weight
classes. For even $m$, the insertion mixture with one flat order. For odd $m$, the uniform
law on the cyclic shifts of $(+3,-3,-1,+1,\dots,-1)$.

\emph{The first odd line.} Only budget $m-1$ is used. By the band at $k=m-1$, the
last-position law satisfies $|\pi_{i}-\frac1m|\le\frac{\eta}{m}$. Hence for any centre $a$,
$\sum_{i}\pi_{i}c_{i}\ge\frac1m\sum_{i}c_{i}-\frac{\eta}{m}\sum_{i}|c_{i}-a|$. With
$c_{i}=|1+g_{i}|$ and $a=2$ the two sums are $m+3$ and $m+1$.

\emph{The second odd line.} Every order pays total excess at least $2+2[\pi_{m}>0]$. The
$+3$ step forces excess $2$ on the two budgets around it, and a positive last item forces two
more at budget $m-1$. The band puts each of the $\frac{m-1}{2}$ positive items last with
probability at least $\frac{1-\eta}{m}$. So the average over budgets, and hence the maximum,
is at least $V^{\mathrm{free}}_{m}+\frac{1-\eta}{mD}$.

\emph{The upper side.} $V_{m}$ is convex in $\eta$, so three laws suffice:
\begin{itemize}
\item the zero-band law at $\eta=0$;
\item the insertion mixture at $\eta=1$, whose value is $V^{\mathrm{free}}_{m}$;
\item a law at the crossing point $\eta_{t}$, with value exactly
$\frac{2(m^{2}+1)}{m^{2}(m-1)D}$.
\end{itemize}
The last law is new. For $m\ge7$ it is a mixture of explicit orders, with masses in units of
$\frac{1}{m^{2}(m-1)}$:
\begin{itemize}
\item mass $m^{2}+1$ on each insertion slot $0,\dots,m-3$ ($m^{2}-m+4$ at slot $0$), split
equally between two versions. Version~A is the insertion order. Version~B is its negation
with the last item set to $-1$;
\item $m(m-1)$ on version~A at slot $m-2$;
\item $2$ on the order that puts the light items of slot $m-2$ first and ends with
$(-3,+3)$;
\item $m-3$ on version~A at slot $0$ with its last two light items swapped, which ends with
$(-1,+1)$.
\end{itemize}
Every budget then carries expected excess exactly $\frac{2(m^{2}+1)}{m^{2}(m-1)}$, and the
class counts of every prefix satisfy the band at $\eta_{t}$. At $m=3$ and $m=5$ an explicit
finite law takes its place.

\section{The flow program has the same value}
\label{sec:flow}

Enumerating orders is hopeless past $m=11$. After symmetrisation, a prefix is characterised,
both for the objective and for the band, by its composition
$v=(n_{-3},n_{-1},n_{+1},n_{+3})\le(1,p,q,1)$. There are $4(p+1)(q+1)=O(m^{2})$ of them. Let
$\mathcal{F}_{m}(\eta)$ be the unit flows on the layered DAG of compositions whose edges add
one item, with the band imposed on the class occupancies.

\begin{theorem}[Flow equivalence]
\label{thm:flow}
Fix $m\ge2$, $\eta$, and a functional $\Phi$ of the budget vector $(\mathbb{E}r_{k})_{k}$. The
infimum of $\Phi$ over band-feasible laws on all $m!$ orders of $g_{m}$ equals its infimum
over $\mathcal{F}_{m}(\eta)$. In particular $V_{m}(\eta)$ is the value of the flow program,
and Theorems~\ref{thm:even} and~\ref{thm:odd} hold for it.
\end{theorem}

\leanbox{\lean{actual\_banded\_infima} in \lean{LeanPaperCFlowBandJoinV1};
\lean{value\_eq\_flow}, \lean{even\_frontier\_flow}, \lean{odd\_breakpoint\_flow} in
\lean{LeanPaperCFlowObjectiveV1}}

Three maps make it exact.
\begin{itemize}
\item A law induces the occupancy of each composition. Conservation is the consistency of the
marginals of one joint law.
\item Averaging over relabellings inside a weight class keeps every $r_{k}$. It replaces each
item's marginal by its class average.
\item A unit flow carries a Markov law on whole trajectories. A trajectory of positive mass is
a sequence of one-item steps, hence the layer map of an actual order. Relabelling within
classes then recovers the per-item band.
\end{itemize}
The two programs reach exactly the same set of budget vectors, and that set is non-empty for
every $\eta\ge0$. The statement concerns the value of the four-type program on $g_{m}$. It
gives no algorithm for perturbed weights.

\begin{table}[t]
\caption{Flow-program sizes against the order formulation. \NUM{}}
\label{tab:agree}
\centering
\footnotesize
\begin{tabular}{@{}rrrl@{}}
\toprule
$m$ & flow states & flow edges & orders \\
\midrule
$11$ & $120$ & $316$ & $13{,}860$ type sequences \\
$15$ & $224$ & $612$ & intractable \\
$21$ & $440$ & $1236$ & intractable \\
$31$ & $960$ & $2756$ & intractable \\
\bottomrule
\end{tabular}
\end{table}

\section{Robustness, and where it ends}
\label{sec:robust}

The family $g_{m}$ is highly structured, and it is fair to ask whether the factor-$m$ price
depends on exact cancellation. Relaxing the constraint is Theorem~\ref{thm:odd}. Perturbing
the weights is this section. The two are not combined.

For arbitrary $g'\in\mathbb{R}^{m}$, let $V(g';\eta)$ be the physical optimum. It
is~\eqref{eq:value} over every order law with the per-item band~\eqref{eq:band}, with the
baseline $b_{k}(g')$ computed from its definition and $D=\|g'\|_{1}$. Perturbations that break
the equal-weight groups are in scope, and $V(g_{m};\eta)=V_{m}(\eta)$
by~\eqref{eq:baseline}. The band is vacuous at $\eta=m$, so $V(g';m)$ is the unconstrained
optimum of $g'$ and $V(g';0)/V(g';m)$ is its price ratio. The normaliser cancels in the
ratio.

\begin{theorem}[Envelope]
\label{thm:robust}
Let $m\ge3$, $g'\in\mathbb{R}^{m}$, $\varepsilon=\|g_{m}-g'\|_{1}$ and
$0<d\le\min(m+4,\|g'\|_{1})$. Then $|V(g';\eta)-V_{m}(\eta)|\le\frac{3\varepsilon}{d}$ for
every $\eta\ge0$. Let $m$ be odd. Whenever $V(g';m)>0$ and $\frac{3\varepsilon}{d}\le V_{m}(0)$,
\[
\frac{V(g';0)}{V(g';m)}\ \ge\ \frac{V_{m}(0)-3\varepsilon/d}{V^{\mathrm{free}}_{m}+3\varepsilon/d},
\]
and for $\varepsilon\le\frac{1}{2m}$, with no further condition,
$\frac{V(g';0)}{V(g';m)}\ge\frac{m-1}{4}$.
\end{theorem}

\leanbox{\lean{family\_value\_robust}, \lean{even\_frontier\_robust} in
\lean{LeanPaperCStabilityFamilyV1}, from \lean{trueValue\_stable} in
\lean{LeanPaperCStabilityEndgameV1}; \lean{physical\_price\_ratio}, \lean{physical\_price\_linear}
in \lean{LeanPaperCRatioJoinV1}}

The constant $3$ pays for three moving parts: the prefix sum, the recomputed baseline and the
normaliser. The envelope alone cannot say where the price collapses, because at fixed
$\varepsilon$ its floor is bounded in $m$. The next theorem locates the collapse from both
sides.

\begin{theorem}[Collapse scale, odd $m$]
\label{thm:collapse}
Let $m\ge3$ be odd.
\begin{enumerate}
\item[(a)] \emph{Upper side.} Take $\delta\in(0,1]$ and move one light item from $-1$ to
$-1+\delta$. This perturbation has $\|g_{m}-g'\|_{1}=\delta$ and a positive unconstrained
value, and its price ratio is at most $\frac{10}{\delta}$.
\item[(b)] \emph{Lower side.} Every $g'$ with $\varepsilon=\|g_{m}-g'\|_{1}\in[\frac1m,\frac16]$
and $V(g';m)>0$ has price ratio at least $\frac{1}{16\varepsilon}$. Every $g'$ with
$\varepsilon\le\frac1m$ and $V(g';m)>0$ has price ratio at least $\frac{m-1}{8}$. So on the
whole range $0<\varepsilon\le\frac16$ the ratio is at least
$\frac1{16}\min\bigl(m-1,\frac1\varepsilon\bigr)$, with no gap in $\varepsilon$.
\end{enumerate}
\end{theorem}

\leanbox{\lean{collapse\_upper}, \lean{collapse\_lower} in \lean{LeanPaperCCollapseScaleV1};
\lean{window\_arith}, \lean{collapse\_window}, \lean{collapse\_all},
\lean{collapse\_window\_satisfiable} in \lean{LeanPaperCCollapseWindowV1}}

Together with the last clause of Theorem~\ref{thm:robust}, this pins the worst price ratio
over perturbations of size $\varepsilon$.
\begin{itemize}
\item For $\varepsilon\le\frac{1}{2m}$ it is at least $\frac{m-1}{4}$.
\item For $\varepsilon\in(\frac{1}{2m},\frac1m)$ it is at least $\frac{m-1}{8}$, over
perturbations with a positive unconstrained value.
\item For $\varepsilon\in[\frac1m,\frac16]$ it lies between $\frac{1}{16\varepsilon}$ and
$\frac{10}{\varepsilon}$. The lower side is over perturbations with a positive unconstrained
value.
\end{itemize}
So the price is linear in $m$ exactly while the perturbation is $O(1/m)$.

The upper side is a parity argument. Every prefix sum of $g_{m}$ has the parity of its
length, which is what makes the free value small. After the move, the best drop of odd size
costs at most $1-\delta$ (drop the moved item), and the best drop of even size still costs
$0$ (leave it out). An order then pays at least $\delta$ at every odd budget before the moved
item and at every even budget from it on. Each pair of budgets $(2i+1,2i+2)$ contains one of
the two. So every order pays at least $\delta\frac{m-1}{2}$ in total, and the free value is at
least $\frac{\delta}{2\|g'\|_{1}}$. The envelope bounds the zero-band value from above.

\paragraph{What a search finds}
We searched for the worst perturbation at $m=5,6,7$, with random restarts and coordinate
descent over the $m!$ formulation and with baseline and normaliser recomputed. A search returns
attained ratios, which are upper bounds on the worst case. At $m=7$ the search keeps $43.5\%$
of the price at $\varepsilon=0.8$ and $30.1\%$ at $\varepsilon=1.6$. At $m=5,6$ no family in
the grid halved it. \NUM{} These sizes lie below the range where the constants of
Theorem~\ref{thm:collapse} separate the two regimes. The search is reported as a consistency
check, not as evidence of the scale.

\section{Choosing the budget after the draw}
\label{sec:expost}

If the budget is chosen after the order is drawn, the regret is
$\mathbb{E}_{P}[\max_{k}r_{k}]$, which is linear in $P$.

\begin{proposition}[Free ex-post value]
\label{prop:expost}
For every $m\ge3$, every order of $g_{m}$ has an interior budget with excess at least $2$,
and some order has excess at most $2$ at every budget. Hence
$\inf_{P}\mathbb{E}_{P}[\max_{0<k<m}r_{k}]=\frac{2}{m+4}=(m-1)V^{\mathrm{free}}_{m}$. For
every law, the ex-ante objective is at most the ex-post one.
\end{proposition}

\leanbox{\lean{pathValue\_eq}, \lean{path\_over\_free}, \lean{exists\_defect},
\lean{objective\_le\_pathObjective} in \lean{LeanPaperCExPostV1}}

\begin{theorem}[Banded ex-post value]
\label{thm:expostband}
Let $V^{\mathrm{post}}_{m}(\eta)$ be the infimum of $\mathbb{E}_{P}[\max_{0<k<m}r_{k}]$ over laws
satisfying~\eqref{eq:band}. For every $\eta\ge0$,
\[
V^{\mathrm{post}}_{m}(\eta)=
\begin{cases}
\dfrac{2}{m+4}, & m\ge4\text{ even},\\[8pt]
\dfrac{2+2(1-\eta)_{+}/m}{m+4}, & m\ge3\text{ odd}.
\end{cases}
\]
\end{theorem}

\leanbox{\lean{expost\_even}, \lean{expost\_odd} in \lean{LeanPaperCExPostBandV1}}

Ex ante, uniformity costs a factor $\Theta(m)$ in odd size. Ex post, the free value is
already $m-1$ times larger, and the band adds at most a factor $1+\frac1m$. So the band price
and the timing price do not multiply.

The lower side is Proposition~\ref{prop:expost} plus one observation. An odd order with the
$+3$ last pays $4$ at budget $m-1$, and the band puts the $+3$ last with probability at least
$\frac{1-\eta}{m}$.

The upper side uses rotations. A uniform law on the $m$ cyclic shifts of one order has exactly
uniform marginals.
\begin{itemize}
\item Every shift of the flat even order $(-3,+1,-1,\dots,+1,-1,+3)$ pays at most $2$ at its
worst budget.
\item Every shift of the odd base order pays at most $2$, except the one that ends with $+3$,
which pays at most $4$.
\item The insertion mixture pays $2$ on every order at $\eta=1$, and blending covers $[0,1]$.
\end{itemize}

\section{Evidence}
\label{sec:evidence}

Every number is produced by a script that checks itself against a statement proved in Lean,
and refuses to write output if a check fails. Figure~\ref{fig:ctheorems} draws the three
closed forms from the theorems. Its markers are floating-point LP solutions at $m\le9$; they
land on the formulas to $10^{-16}$ and certify nothing. Figure~\ref{fig:cstab} checks the
envelope against every stored perturbation.

\begin{table}[t]
\caption{Evidence packages. Each script self-checks against a Lean statement.}
\label{tab:evidence}
\centering
\footnotesize
\resizebox{\columnwidth}{!}{%
\begin{tabular}{@{}llr@{}}
\toprule
Package & Script & Checks \\
\midrule
Formula cross-check (ex ante, ex post) & \texttt{EXP\_paperC\_formula\_check} & 2/2 \\
Theorem figure & \texttt{FIG\_paperC\_theorems} & 3/3 \\
Flow program, breakpoints & \texttt{EXP\_paperC\_flow\_lp} & 5/5 \\
Breakpoint certificates & \texttt{EXP\_paperC\_breakpoint\_certs} & 5/5 \\
Exact Farkas roots & \texttt{EXP\_paperC\_exact\_dual} & 5/5 \\
Stored data vs.\ theorems & \texttt{FIG\_paperC\_endpoints} & 10/10 \\
Perturbation, both normalisations & \texttt{EXP\_paperC\_perturbation} & 9/9 \\
Adversarial price search & \texttt{EXP\_paperC\_robustness} & 4/4 \\
Window sums, value & \texttt{EXP\_paperC\_window} & 7/7 \\
\bottomrule
\end{tabular}}
\end{table}

\begin{figure*}[t]
\centering
\includegraphics[width=0.95\textwidth]{paperC_theorems.pdf}
\caption{The three closed forms, drawn from the theorems. Left: the whole frontier
$V_{m}(\eta)/V^{\mathrm{free}}_{m}$ (Theorems~\ref{thm:even}, \ref{thm:odd}), odd sizes solid
and even sizes dashed. Centre: the banded ex-post value over the free ex-post value
$\frac{2}{m+4}$ (Theorem~\ref{thm:expostband}); it is at most $1+\frac1m$, and exactly $1$ for
even $m$. Markers: floating-point LP solutions at $m\le9$ (\NUM), cross-checks only. Right:
the collapse scale at $m=101$, theorem curves only. Black: the envelope floor of
Theorem~\ref{thm:robust} with $d=m+4-\varepsilon$. Green: the linear guarantee $\frac{m-1}{4}$
for $\varepsilon\le\frac{1}{2m}$. Blue: the lower side $\frac{1}{16\varepsilon}$ on
$[\frac1m,\frac16]$. Dashed: the one-item construction $\frac{10}{\varepsilon}$. Dotted: the
unperturbed ratio. (\texttt{FIG\_paperC\_theorems}, 3/3.)}
\label{fig:ctheorems}
\end{figure*}

\begin{figure*}[t]
\centering
\includegraphics[width=0.92\textwidth]{paperC_stability.pdf}
\caption{The envelope against stored perturbations; nothing was re-solved. Left: all 220
perturbations of the physical objective, 160 group perturbations over the flow LP and 60
per-item perturbations over the $m!$ order LP. They are drawn against $3\|g_{m}-g'\|_{1}/d$
with $d=\min(\|g_{m}\|_{1},\|g'\|_{1})$. A point in the shaded region would refute
Theorem~\ref{thm:robust}; none enters it, and the worst ratio is $0.32$. The producer skips
solver failures without recording them, so row counts were compared with its fixed grid, and
none was dropped. Centre: the same drifts divided by the envelope. Objectives with the
normaliser or the baseline frozen are drawn only to show which object moves. Right: price
ratios attained by the adversarial search (markers), which are upper bounds on the worst case,
against the fixed-normaliser floor $\frac{1}{4\varepsilon+4/(m-1)}$ (lines, odd $m$; its
status is stated in the supplement).}
\label{fig:cstab}
\end{figure*}

\section{Limitations}
\label{sec:limits}

\begin{enumerate}
\item[1)] \textbf{One family.} Every exact value is for $g_{m}$ of~\eqref{eq:family}. Other
weights enter only through Theorems~\ref{thm:robust} and~\ref{thm:collapse}.
\item[2)] \textbf{Constants of the collapse.} Theorem~\ref{thm:collapse} pins the collapse
scale in order. Its two sides differ by a constant factor, and its lower side needs a
positive unconstrained value, which the unperturbed family has.
\item[3)] \textbf{Value, not algorithm.} Theorem~\ref{thm:flow} equates two values on
$g_{m}$. It gives no polynomial-size program for perturbed weights.
\item[4)] \textbf{Two objectives, kept apart.} Every ex-ante statement concerns
$\max_{k}\mathbb{E}r_{k}$, and every ex-post one concerns $\mathbb{E}\max_{k}r_{k}$.
\end{enumerate}

\section{Conclusion}

Suppose every item must be dropped about equally often at every budget. That costs a
bounded factor in even size and a linear factor in odd size, and the whole frontier in the
band width is known in closed form at every size. The value over all $m!$ orders is the value
of an $O(m^{2})$ flow. The price survives perturbations of the weights exactly while they are
$O(1/m)$, and it decays as $\Theta(1/\varepsilon)$ beyond. When the budget is chosen after
the draw, the band is almost free, so the timing price and the band price do not multiply.

\section*{Companion material}
The supplement gives the full proof architecture of each theorem, every supporting Lean
declaration, and the frozen builds it was checked in. The claim map records, for each
headline, its model, its declaration, its scope and its evidence. Corrections made during
preparation are recorded in the project ledger, not in this text.

\begin{thebibliography}{00}

\bibitem{singh2018} A. Singh and T. Joachims, ``Fairness of exposure in rankings,'' in
\emph{Proc. 24th ACM SIGKDD Int. Conf. Knowledge Discovery and Data Mining}, 2018,
pp.~2219--2228.

\bibitem{celis2018} L. E. Celis, D. Straszak, and N. K. Vishnoi, ``Ranking with fairness
constraints,'' in \emph{Proc. 45th Int. Colloq. Automata, Languages, and Programming
(ICALP)}, 2018.

\bibitem{singh2021} A. Singh, D. Kempe, and T. Joachims, ``Fairness in ranking under
uncertainty,'' in \emph{Advances in Neural Information Processing Systems 34}, 2021,
pp.~11896--11908.

\bibitem{freeman2020} R. Freeman, N. Shah, and R. Vaish, ``Best of both worlds: Ex-ante and
ex-post fairness in resource allocation,'' in \emph{Proc. 21st ACM Conf. Economics and
Computation (EC)}, 2020.

\bibitem{schrijver2003} A. Schrijver, \emph{Combinatorial Optimization: Polyhedra and
Efficiency}. Berlin, Germany: Springer, 2003.

\bibitem{mathlib2020} The mathlib Community, ``The Lean mathematical library,'' in
\emph{Proc. 9th ACM SIGPLAN Int. Conf. Certified Programs and Proofs (CPP)}, 2020,
pp.~367--381.

\bibitem{demoura2021} L. de Moura and S. Ullrich, ``The Lean 4 theorem prover and
programming language,'' in \emph{Proc. 28th Int. Conf. Automated Deduction (CADE)}, 2021,
pp.~625--635.

\end{thebibliography}

\end{document}
